\section{Cross-episode experience and persistent memory}
\label{sec:fly}
Foundation language models are stateless across sessions: once an interaction ends,
the model retains no memory of whether its previous trading strategy made or lost money.
Fine-tuning the entire LLM with gradient descent after each trading day is impractical,
computationally expensive, and prone to overfitting on noisy financial markets.

To provide cross-episode memory without altering LLM weights, we use an external
associative memory module inspired by the fruit-fly mushroom body \cite{huang2024memory}.
In biological systems, the mushroom body maps complex sensory features into sparse
representations and uses dopamine-gated synaptic plasticity to associate sensory cues
with delayed rewards and punishments across multiple decay rates (short- and long-term
memory).

\paragraph{Associative memory architecture.}
Our memory module projects market state features and candidate actions into a high-dimensional
sparse representation. When the agent evaluates a new trading decision, it queries the
memory module for associative preference scores based on prior experiences in similar
market regimes. These scores serve as contextual guidance alongside the prompt, without
modifying the frozen LLM backbone.

\paragraph{Updating from complete holding episodes.}
Financial rewards are inherently delayed: an allocation chosen today only yields returns
after a multi-day holding period. The feedback interface models an entry and its subsequent
holding days as a single complete decision episode. For positive account values
$W_0,\ldots,W_T$, the target return is
\begin{equation}
G = \sum_{t=0}^{T-1}\log(W_{t+1}/W_t)
  = \log(W_T/W_0).
\end{equation}
This return accounts for transaction fees, daily price fluctuations, and final liquidation.
Memory updates fire only after the position is fully closed and settlement clears.

The verification interface checks timestamp ordering, return accounting, and data
lineage. Trades relying on look-ahead data or unadjusted prices are rejected before
reaching the memory module, preventing spurious profits from corrupting associative
weights. Valid negative returns remain active learning signals, allowing the agent to
learn from losses. Each verified trade receipt updates memory exactly once.

\paragraph{Isolating learning from risk filtering.}
To verify that performance gains reflect genuine associative learning rather than
passive risk reduction, our benchmark compares the updating memory against a frozen-memory
control and a volatility exposure gate under identical historical splits and costs.
This setup isolates experiential adaptation from static volatility filtering.
